
\Data\ID{1003138501}
\Updated{2020-07-15 03:07:21}
\Level{A}
\SubArea{Logarithmic equations and inequalities}
\LANGen{
\Question{
Solve.
\[ \log_2(3x-5)=4\]}
{\( x=7 \)}{\( x=3 \)}{\( x=\frac{11}3 \)}{\( x=-\frac13 \)}{}{}}
\LANGpl{
\Question{
Rozwiąż.
\[ \log_2(3x-5)=4 \]}
{\( x=7 \)}{\( x=3 \)}{\( x=\frac{11}3 \)}{\( x=-\frac13 \)}{}{}}
\LANGcs{
\Question{
Řešte danou rovnici
\[ \log_2(3x-5)=4 \]}
{\( x=7 \)}{\( x=3 \)}{\( x=\frac{11}3 \)}{\( x=-\frac13 \)}{}{}}
\LANGsk{
\Question{
Riešte danú rovnicu.
\[ \log_2(3x-5)=4\]}
{\( x=7 \)}{\( x=3 \)}{\( x=\frac{11}3 \)}{\( x=-\frac13 \)}{}{}}
\LANGes{
\Question{
Resuelve.
\[ \log_2(3x-5)=4\]}
{\( x=7 \)}{\( x=3 \)}{\( x=\frac{11}3 \)}{\( x=-\frac13 \)}{}{}}
\End

\Data\ID{1003138502}
\Updated{2020-07-15 03:07:21}
\Level{A}
\SubArea{Logarithmic equations and inequalities}
\LANGen{
\Question{
Solve.
\[ \log_{\frac13}\!(3-x)=0\]}
{\( x=2 \)}{The equation has no solution.}{\( x =3 \)}{\( x=-4 \)}{}{}}
\LANGpl{
\Question{
Rozwiąż.
\[ \log_{\frac13}\!(3-x)=0 \]}
{\( x=2 \)}{Równanie nie ma rozwiązania.}{\( x =3 \)}{\( x=-4 \)}{}{}}
\LANGcs{
\Question{
Řešte danou rovnici.
\[ \log_{\frac13}\!(3-x)=0 \]}
{\( x=2 \)}{Rovnice nemá řešení.}{\( x =3 \)}{\( x=-4 \)}{}{}}
\LANGsk{
\Question{
Riešte danú rovnicu.
\[ \log_{\frac13}\!(3-x)=0 \]}
{\( x=2 \)}{Rovnica nemá riešenie.}{\( x =3 \)}{\( x=-4 \)}{}{}}
\LANGes{
\Question{
Resuelve. 
\[ \log_{\frac13}\!(3-x)=0\]}
{\( x=2 \)}{La ecuación no tiene solución.}{\( x =3 \)}{\( x=-4 \)}{}{}}
\End

\Data\ID{1003138503}
\Updated{2020-07-15 03:07:21}
\Level{A}
\SubArea{Logarithmic equations and inequalities}
\LANGen{
\Question{
How many solutions in the set of integers does the following equation have?
\[ \log_{\frac12}\!(x+6)=\log_{\frac12}\!(3x-6) \]}
{exactly one positive solution}{exactly one negative solution}{exactly one solution equal to zero}{no solutions}{}{}}
\LANGpl{
\Question{
Ile rozwiązań w zbiorze liczb całkowitych posiada podane równanie?
\[ \log_{\frac12}\!(x+6)=\log_{\frac12}\!(3x-6) \]}
{dokładnie jedno dodatnie rozwiązanie}{dokładnie jedno ujemne rozwiązanie}{dokładnie jedno rozwiązanie równe zeru}{nie ma rozwiązania}{}{}}
\LANGcs{
\Question{
Kolik řešení v množině celých čísel má následující rovnice?
\[ \log_{\frac12}\!(x+6)=\log_{\frac12}\!(3x-6) \]}
{právě jedno řešení}{právě jedno záporné řešení}{právě jedno nulové řešení}{žádné řešení}{}{}}
\LANGsk{
\Question{
Koľko riešení na množine celých čísel má nasledujúca rovnica?
\[ \log_{\frac12}\!(x+6)=\log_{\frac12}\!(3x-6) \]}
{práve jedno riešenie}{práve jedno záporné riešenie}{práve jedno nulové riešenie}{žiadne riešenie}{}{}}
\LANGes{
\Question{
¿Cuántas soluciones en el conjunto de los números enteros tiene la siguiente ecuación? \[ \log_{\frac12}\!(x+6)=\log_{\frac12}\!(3x-6) \]}
{exactamente una solución positiva}{exactamente una solución negativa}{exactamente una solución igual a cero}{ninguna solución}{}{}}
\End

\Data\ID{1003138504}
\Updated{2020-07-15 03:07:21}
\Level{A}
\SubArea{Logarithmic equations and inequalities}
\LANGen{
\Question{
Solve.
\[ \log_3(-5x-2)=\log_3(x-4) \]}
{The equation has no solution.}{\( x=\frac13 \)}{\( x=3 \)}{\( x=\frac32 \)}{}{}}
\LANGpl{
\Question{
Rozwiąż.
\[ \log_3(-5x-2)=\log_3(x-4)\]}
{Równanie nie ma rozwiązania.}{\( x=\frac13 \)}{\( x=3 \)}{\( x=\frac32 \)}{}{}}
\LANGcs{
\Question{
Řešte danou rovnici.
\[ \log_3(-5x-2)=\log_3(x-4) \]}
{Rovnice nemá řešení.}{\( x=\frac13 \)}{\( x=3 \)}{\( x=\frac32 \)}{}{}}
\LANGsk{
\Question{
Riešte danú rovnicu.
\[ \log_3(-5x-2)=\log_3(x-4)\]}
{Rovnica nemá riešenie.}{\( x=\frac13 \)}{\( x=3 \)}{\( x=\frac32 \)}{}{}}
\LANGes{
\Question{
Resuelve.
\[ \log_3(-5x-2)=\log_3(x-4) \]}
{La ecuación no tiene solución.}{\( x=\frac13 \)}{\( x=3 \)}{\( x=\frac32 \)}{}{}}
\End

\Data\ID{1003138505}
\Updated{2020-07-15 03:07:21}
\Level{A}
\SubArea{Logarithmic equations and inequalities}
\LANGen{
\Question{
Solve.
\[ \ln(-2x-6)=\ln(-3x-9) \]}
{The equation has no solution.}{\( x=-3 \)}{\( x=3 \)}{\( x=-15 \)}{}{}}
\LANGpl{
\Question{
Rozwiąż.
\[ \ln(-2x-6)=\ln(-3x-9) \]}
{Równanie nie ma rozwiązania.}{\( x=-3 \)}{\( x=3 \)}{\( x=-15 \)}{}{}}
\LANGcs{
\Question{
Řešte danou rovnici.
\[ \ln(-2x-6)=\ln(-3x-9) \]}
{Rovnice nemá řešení.}{\( x=-3 \)}{\( x=3 \)}{\( x=-15 \)}{}{}}
\LANGsk{
\Question{
Riešte danú rovnicu.
\[ \ln(-2x-6)=\ln(-3x-9) \]}
{Rovnica nemá riešenie.}{\( x=-3 \)}{\( x=3 \)}{\( x=-15 \)}{}{}}
\LANGes{
\Question{
Resuelve.
\[ \ln(-2x-6)=\ln(-3x-9) \]}
{La ecuación no tiene solución.}{\( x=-3 \)}{\( x=3 \)}{\( x=-15 \)}{}{}}
\End

\Data\ID{1003138506}
\Updated{2020-07-15 03:07:21}
\Level{A}
\SubArea{Logarithmic equations and inequalities}
\LANGen{
\Question{
Solve.
\[ \frac{\log_3x-3}{3+\log_3x}=0\]}
{\( x=27 \)}{\( x=9 \)}{\( x_1=\frac1{27};\ x_2=27 \)}{infinitely many solutions}{}{}}
\LANGpl{
\Question{
Rozwiąż.
\[ \frac{\log_3x-3}{3+\log_3x}=0\]}
{\( x=27 \)}{\( x=9 \)}{\( x_1=\frac1{27};\ x_2=27 \)}{nieskończenie wiele rozwiązań}{}{}}
\LANGcs{
\Question{
Řešte danou rovnici.
\[ \frac{\log_3x-3}{3+\log_3x}=0 \]}
{\( x=27 \)}{\( x=9 \)}{\( x_1=\frac1{27};\ x_2=27 \)}{Rovnice má nekonečně mnoho řešení.}{}{}}
\LANGsk{
\Question{
Riešte danú rovnicu.
\[ \frac{\log_3x-3}{3+\log_3x}=0 \]}
{\( x=27 \)}{\( x=9 \)}{\( x_1=\frac1{27};\ x_2=27 \)}{Rovnica má nekonečne veľa riešení.}{}{}}
\LANGes{
\Question{
Resuelve.
\[ \frac{\log_3x-3}{3+\log_3x}=0\]}
{\( x=27 \)}{\( x=9 \)}{\( x_1=\frac1{27};\ x_2=27 \)}{Soluciones infinitas.}{}{}}
\End

\Data\ID{1003138507}
\Updated{2020-07-15 03:07:21}
\Level{A}
\SubArea{Logarithmic equations and inequalities}
\LANGen{
\Question{
Solve.
\[ \log_4x=2-\log_48\]}
{\( x=2 \)}{\( x=1 \)}{\( x=0 \)}{The equation has no solution.}{}{}}
\LANGpl{
\Question{
Rozwiąż.
\[ \log_4x=2-\log_48\]}
{\( x=2 \)}{\( x=1 \)}{\( x=0 \)}{Równanie nie ma rozwiązania.}{}{}}
\LANGcs{
\Question{
Řešte danou rovnici.
\[ \log_4x=2-\log_48\]}
{\( x=2 \)}{\( x=1 \)}{\( x=0 \)}{Rovnice nemá řešení.}{}{}}
\LANGsk{
\Question{
Riešte danú rovnicu.
\[ \log_4x=2-\log_48\]}
{\( x=2 \)}{\( x=1 \)}{\( x=0 \)}{Rovnica nemá riešenie.}{}{}}
\LANGes{
\Question{
Resuelve.
\[ \log_4x=2-\log_48\]}
{\( x=2 \)}{\( x=1 \)}{\( x=0 \)}{La ecuación no tiene solución.}{}{}}
\End

\Data\ID{1003138508}
\Updated{2020-07-15 03:07:21}
\Level{A}
\SubArea{Logarithmic equations and inequalities}
\LANGen{
\Question{
How many solutions does the following equation have?
\[ 2\log x^4-\log x^5-3\log x^3= 12 \]}
{exactly one positive solution}{exactly one negative solution}{exactly one integer solution}{no solutions}{}{}}
\LANGpl{
\Question{
Ile rozwiązań ma podane równanie?
\[ 2\log x^4-\log x^5-3\log x^3= 12 \]}
{dokładnie jedno dodatnie rozwiązanie}{dokładnie jedno ujemne rozwiązanie}{rozwiązaniem jest jedna liczba całkowita}{brak rozwiązania}{}{}}
\LANGcs{
\Question{
Kolik řešení má daná rovnice?
\[ 2\log x^4-\log x^5-3\log x^3= 12 \]}
{právě jedno kladné řešení}{právě jedno záporné řešení}{právě jedno celočíselné řešení}{nemá řešení}{}{}}
\LANGsk{
\Question{
Koľko riešení má daná rovnica?
\[ 2\log x^4-\log x^5-3\log x^3= 12 \]}
{práve jedno kladné riešenie}{práve jedno záporné riešenie}{práve jedno celočíselné riešenie}{nemá riešenie}{}{}}
\LANGes{
\Question{
¿Cuántas soluciones tiene la siguiente ecuación? 
\[ 2\log x^4-\log x^5-3\log x^3= 12 \]}
{exactamente una solución positiva}{exactamente una solución negativa}{exactamente una solución del conjunto de los números enteros}{sin solución}{}{}}
\End

\Data\ID{1003138509}
\Updated{2020-07-15 03:07:21}
\Level{A}
\SubArea{Logarithmic equations and inequalities}
\LANGen{
\Question{
Solve.
\[ \log_2(x-1)+\log_2x=1 \]}
{\( x=2 \)}{\( x_1=-1;\ x_2=2 \)}{\( x_1=-2;\ x_2=1 \)}{\( x=1 \)}{}{}}
\LANGpl{
\Question{
Rozwiąż.
\[ \log_2(x-1)+\log_2x=1\]}
{\( x=2 \)}{\( x_1=-1;\ x_2=2 \)}{\( x_1=-2;\ x_2=1 \)}{\( x=1 \)}{}{}}
\LANGcs{
\Question{
Řešte danou rovnici.
\[ \log_2(x-1)+\log_2x=1 \]}
{\( x=2 \)}{\( x_1=-1;\ x_2=2 \)}{\( x_1=-2;\ x_2=1 \)}{\( x=1 \)}{}{}}
\LANGsk{
\Question{
Riešte danú rovnicu.
\[ \log_2(x-1)+\log_2x=1 \]}
{\( x=2 \)}{\( x_1=-1;\ x_2=2 \)}{\( x_1=-2;\ x_2=1 \)}{\( x=1 \)}{}{}}
\LANGes{
\Question{
Resuelve.
\[ \log_2(x-1)+\log_2x=1 \]}
{\( x=2 \)}{\( x_1=-1;\ x_2=2 \)}{\( x_1=-2;\ x_2=1 \)}{\( x=1 \)}{}{}}
\End

\Data\ID{1003138510}
\Updated{2020-07-15 03:07:21}
\Level{A}
\SubArea{Logarithmic equations and inequalities}
\LANGen{
\Question{
How many solutions does the following equation have?
\[ \log(x-2)+\log(x+2)=2\log(2-x) \]}
{no solutions}{exactly one solution}{exactly two solutions}{infinitely many solutions}{}{}}
\LANGpl{
\Question{
Ile rozwiązań ma poniższe równanie?
\[ \log(x-2)+\log(x+2)=2\log(2-x) \]}
{nie ma rozwiązania}{dokładnie jedno rozwiązanie}{dokładnie dwa rozwiązania}{nieskończenie wiele rozwiązań}{}{}}
\LANGcs{
\Question{
Kolik řešení má daná rovnice?
\[ \log(x-2)+\log(x+2)=2\log(2-x) \]}
{nemá řešení}{právě jedno řešení}{právě dvě řešení}{nekonečně mnoho řešení}{}{}}
\LANGsk{
\Question{
Koľko riešení má daná rovnica?
\[ \log(x-2)+\log(x+2)=2\log(2-x) \]}
{nemá riešenie}{práve jedno riešenie}{práve dve riešenia}{nekonečne veľa riešení}{}{}}
\LANGes{
\Question{
¿Cuántas soluciones tiene la siguiente ecuación? 
\[ \log(x-2)+\log(x+2)=2\log(2-x) \]}
{sin solución}{exactamente una solución}{exactamente dos soluciones}{soluciones infinitas}{}{}}
\End

\Data\ID{1003138512}
\Updated{2020-07-15 03:07:21}
\Level{A}
\SubArea{Logarithmic equations and inequalities}
\LANGen{
\Question{
How many solutions in the set of integers does the following equation have?
\[ \log_3(x^2-4x)-1=\log_3(2-x) \]}
{exactly one negative solution}{exactly one positive solution}{exactly two solutions}{no solutions}{}{}}
\LANGpl{
\Question{
Ile rozwiązań w zbiorze liczb całkowitych posiada poniższe równanie? 
\[ \log_3(x^2-4x)-1=\log_3(2-x) \]}
{dokładnie jedno ujemne rozwiązanie}{dokładnie jedno pozytywne rozwiązanie}{dokładnie dwa rozwiązania}{nie ma rozwiązania}{}{}}
\LANGcs{
\Question{
Kolik řešení na množině celých čísel má následující rovnice?
\[ \log_3(x^2-4x)-1=\log_3(2-x) \]}
{právě jedno záporné řešení}{právě jedno kladné řešení}{právě dvě řešení}{nemá řešení}{}{}}
\LANGsk{
\Question{
Koľko riešení na množine celých čísel má nasledujúca rovnica?
\[ \log_3(x^2-4x)-1=\log_3(2-x) \]}
{práve jedno záporné riešenie}{práve jedno kladné riešenie}{práve dve riešenia}{nemá riešenie}{}{}}
\LANGes{
\Question{
¿Cuántas soluciones en el conjunto de los números enteros tiene la siguiente ecuación? 
\[ \log_3(x^2-4x)-1=\log_3(2-x) \]}
{Exactamente una solución negativa.}{Exactamente una solución positiva.}{Exactamente dos soluciones.}{Sin solución.}{}{}}
\End

\Data\ID{1003138513}
\Updated{2020-07-15 03:07:21}
\Level{A}
\SubArea{Logarithmic equations and inequalities}
\LANGen{
\Question{
How many of the following equations have the solution that is a prime number?
\[ \begin{aligned}
\log_2\!\left(3^x-1\right)&=3 \\
\log_3\!\left(8^x+1\right)&=2 \\
\log_{\frac12}\!\left(2^x+8\right)&=-4 \end{aligned} \]}
{\( 2 \)}{\( 3 \)}{\( 1 \)}{\( 0 \)}{}{}}
\LANGpl{
\Question{
Dla ilu z poniższych równań rozwiązaniem jest liczba pierwsza?
\[ \begin{aligned}
\log_2\!\left(3^x-1\right)&=3 \\
\log_3\!\left(8^x+1\right)&=2 \\
\log_{\frac12}\!\left(2^x+8\right)&=-4 \end{aligned} \]}
{\( 2 \)}{\( 3 \)}{\( 1 \)}{\( 0 \)}{}{}}
\LANGcs{
\Question{
Kolik z daných rovnic má řešení, které je prvočíslem?
\[ \begin{aligned}
\log_2\!\left(3^x-1\right)&=3 \\
\log_3\!\left(8^x+1\right)&=2 \\
\log_{\frac12}\!\left(2^x+8\right)&=-4 \end{aligned} \]}
{\( 2 \)}{\( 3 \)}{\( 1 \)}{\( 0 \)}{}{}}
\LANGsk{
\Question{
Koľko z daných rovníc má riešenie prvočíslo?
\[ \begin{aligned}
\log_2\!\left(3^x-1\right)&=3 \\
\log_3\!\left(8^x+1\right)&=2 \\
\log_{\frac12}\!\left(2^x+8\right)&=-4 \end{aligned} \]}
{\( 2 \)}{\( 3 \)}{\( 1 \)}{\( 0 \)}{}{}}
\LANGes{
\Question{
¿Cuántas de las siguientes ecuaciones tienen la solución que sea número primo?
\[ \begin{aligned}
\log_2\!\left(3^x-1\right)&=3 \\
\log_3\!\left(8^x+1\right)&=2 \\
\log_{\frac12}\!\left(2^x+8\right)&=-4 \end{aligned} \]}
{\( 2 \)}{\( 3 \)}{\( 1 \)}{\( 0 \)}{}{}}
\End

\Data\ID{1003138514}
\Updated{2020-10-19 04:10:14}
\Level{A}
\SubArea{Logarithmic equations and inequalities}
\LANGen{
\Question{
How many solutions in the set of integers does the following equation have?
\[ \log_3\frac{2x^2+4}{x^2-4}=1 \]}
{exactly two solutions}{exactly one solution}{no solutions}{exactly one solution equal to zero}{}{}}
\LANGpl{
\Question{
Ile rozwiązań w zbiorze liczb całkowitych posiada poniższe równanie?
\[ \log_3\frac{2x^2+4}{x^2-4}=1 \]}
{dokładnie dwa rozwiązania}{dokładnie jedno rozwiązanie}{nie ma rozwiązania}{dokładnie jedno rozwiązanie równe zeru}{}{}}
\LANGcs{
\Question{
Kolik řešení na množině celých čísel má následující rovnice?
\[ \log_3\frac{2x^2+4}{x^2-4}=1 \]}
{právě dvě řešení}{právě jedno řešení}{nemá řešení}{právě jedno nulové řešení}{}{}}
\LANGsk{
\Question{
Koľko riešení na množine celých čísel má nasledujúca rovnica?
\[ \log_3\frac{2x^2+4}{x^2-4}=1 \]}
{práve dve riešenia}{práve jedno riešenie}{nemá riešenie}{práve jedno nulové riešenie}{}{}}
\LANGes{
\Question{
¿Cuántas soluciones en el conjunto de los números enteros tiene la siguiente ecuación? \[ \log_3\frac{2x^2+4}{x^2-4}=1 \]}
{Exactamente dos soluciones.}{Exactamente una solución.}{Sin solución.}{Exactamente una solución igual a cero.}{}{}}
\End

\Data\ID{1003138511}
\Updated{2020-07-15 03:07:21}
\Level{B}
\SubArea{Logarithmic equations and inequalities}
\LANGen{
\Question{
Solve
\[ \frac{\log(4x)}{\log(x-3)}=\frac{\log16}{\log4} \text{ .} \]}
{\( x = 9 \)}{\( x_1=1;\ x_2=9 \)}{\( x_1=-9;\ x_2=-1 \)}{The equation has no solution.}{}{}}
\LANGpl{
\Question{
Rozwiąż
\[ \frac{\log(4x)}{\log(x-3)}=\frac{\log16}{\log4} \text{ .} \]}
{\( x = 9 \)}{\( x_1=1;\ x_2=9 \)}{\( x_1=-9;\ x_2=-1 \)}{Równanie nie ma rozwiązania.}{}{}}
\LANGcs{
\Question{
Řešte rovnici
\[ \frac{\log(4x)}{\log(x-3)}=\frac{\log16}{\log4} \text{ .} \]}
{\( x = 9 \)}{\( x_1=1;\ x_2=9 \)}{\( x_1=-9;\ x_2=-1 \)}{Rovnice nemá řešení.}{}{}}
\LANGsk{
\Question{
Riešte rovnicu
\[ \frac{\log(4x)}{\log(x-3)}=\frac{\log16}{\log4} \text{ .} \]}
{\( x = 9 \)}{\( x_1=1;\ x_2=9 \)}{\( x_1=-9;\ x_2=-1 \)}{Rovnica nemá riešenie.}{}{}}
\LANGes{
\Question{
Resuelve.
\[ \frac{\log(4x)}{\log(x-3)}=\frac{\log16}{\log4} \text{ .} \]}
{\( x = 9 \)}{\( x_1=1;\ x_2=9 \)}{\( x_1=-9;\ x_2=-1 \)}{La ecuación no tiene solución.}{}{}}
\End

\Data\ID{1003143201}
\Updated{2020-07-15 03:07:21}
\Level{B}
\SubArea{Logarithmic equations and inequalities}
\LANGen{
\Question{
Solve
\[ \log_3\!\left(3\log_3x+6\right)=1 \text{ .}\]}
{\( x=\frac13 \)}{\( x=27 \)}{\( x=9 \)}{\( x=\frac19 \)}{}{}}
\LANGpl{
\Question{
Rozwiąż
\[ \log_3\!\left(3\log_3x+6\right)=1 \text{ .}\]}
{\( x=\frac13 \)}{\( x=27 \)}{\( x=9 \)}{\( x=\frac19 \)}{}{}}
\LANGcs{
\Question{
Řešte rovnici
\[ \log_3\!\left(3\log_3x+6\right)=1 \text{ .}\]}
{\( x=\frac13 \)}{\( x=27 \)}{\( x=9 \)}{\( x=\frac19 \)}{}{}}
\LANGsk{
\Question{
Riešte rovnicu
\[ \log_3\!\left(3\log_3x+6\right)=1 \text{ .}\]}
{\( x=\frac13 \)}{\( x=27 \)}{\( x=9 \)}{\( x=\frac19 \)}{}{}}
\LANGes{
\Question{
Resuelve.
\[ \log_3\!\left(3\log_3x+6\right)=1 \text{ .}\]}
{\( x=\frac13 \)}{\( x=27 \)}{\( x=9 \)}{\( x=\frac19 \)}{}{}}
\End

\Data\ID{1003143202}
\Updated{2020-07-15 03:07:21}
\Level{B}
\SubArea{Logarithmic equations and inequalities}
\LANGen{
\Question{
What is the total count of all integer solutions of the all following equations? 
\[ \begin{aligned}
\log_{\frac12}\!\left(\log_2x\right)&=-1 \\
\log_5\!\left(\log_{\frac15}x\right)&=0 \\
-\log_{\frac13}\!\left(\log_{\frac12}x\right)&=1 \end{aligned}\]}
{\( 1 \)}{\( 2 \)}{\( 3 \)}{\( 0 \)}{}{}}
\LANGpl{
\Question{
Jaka jest ostateczna liczba wszystkich liczb całkowitych będących rozwiązaniem podanego równania?  
\[ \begin{aligned}
\log_{\frac12}\!\left(\log_2x\right)&=-1 \\
\log_5\!\left(\log_{\frac15}x\right)&=0 \\
-\log_{\frac13}\!\left(\log_{\frac12}x\right)&=1 \end{aligned}\]}
{\( 1 \)}{\( 2 \)}{\( 3 \)}{\( 0 \)}{}{}}
\LANGcs{
\Question{
Určete celkový počet celočíselných řešení daných rovnic.
\[ \begin{aligned}
\log_{\frac12}\!\left(\log_2x\right)&=-1 \\
\log_5\!\left(\log_{\frac15}x\right)&=0 \\
-\log_{\frac13}\!\left(\log_{\frac12}x\right)&=1 \end{aligned}\]}
{\( 1 \)}{\( 2 \)}{\( 3 \)}{\( 0 \)}{}{}}
\LANGsk{
\Question{
Určte celkový počet celočíselných riešení daných rovníc.
\[ \begin{aligned}
\log_{\frac12}\!\left(\log_2x\right)&=-1 \\
\log_5\!\left(\log_{\frac15}x\right)&=0 \\
-\log_{\frac13}\!\left(\log_{\frac12}x\right)&=1 \end{aligned}\]}
{\( 1 \)}{\( 2 \)}{\( 3 \)}{\( 0 \)}{}{}}
\LANGes{
\Question{
Identifica el número total de las soluciones en números enteros de todas las ecuaciones siguientes.  
\[ \begin{aligned}
\log_{\frac12}\!\left(\log_2x\right)&=-1 \\
\log_5\!\left(\log_{\frac15}x\right)&=0 \\
-\log_{\frac13}\!\left(\log_{\frac12}x\right)&=1 \end{aligned}\]}
{\( 1 \)}{\( 2 \)}{\( 3 \)}{\( 0 \)}{}{}}
\End

\Data\ID{1003143203}
\Updated{2020-07-15 03:07:21}
\Level{B}
\SubArea{Logarithmic equations and inequalities}
\LANGen{
\Question{
Which of the following statements is not true about the equation?
\[ \log_32=\log_3\!\left(3-\log_2x\right) \]}
{The solution is an odd number.}{The solution is \( x=2 \).}{The solution is an even number.}{The solution is a prime number.}{}{}}
\LANGpl{
\Question{
Które z poniższych stwierdzeń nie jest prawdziwe w odniesieniu do poniższego równania?
\[ \log_32=\log_3\!\left(3-\log_2x\right) \]}
{Rozwiązaniem jest liczba nieparzysta.}{Rozwiązaniem jest  \( x=2 \).}{Rozwiązaniem jest liczba parzysta.}{Rozwiązaniem jest liczba pierwsza.}{}{}}
\LANGcs{
\Question{
Které z následujících tvrzení o rovnici není pravdivé?
\[ \log_32=\log_3\!\left(3-\log_2x\right) \]}
{Řešením rovnice je liché číslo.}{Řešením rovnice je \( x=2 \).}{Řešením rovnice je sudé číslo.}{Řešením rovnice je prvočíslo.}{}{}}
\LANGsk{
\Question{
Ktoré z nasledujúcich tvrdení o rovnici nie je pravdivé?
\[ \log_32=\log_3\!\left(3-\log_2x\right) \]}
{Riešením rovnice je nepárne číslo.}{Riešením rovnice je \( x=2 \).}{Riešením rovnice je párne číslo.}{Riešením rovnice je prvočíslo.}{}{}}
\LANGes{
\Question{
¿Cuál de los enunciados de la siguiente ecuación no es verdadero?
\[ \log_32=\log_3\!\left(3-\log_2x\right) \]}
{La solución en un número impar.}{La solución es \( x=2 \).}{La solución es un número par.}{La solución es un número primo.}{}{}}
\End

\Data\ID{1003158801}
\Updated{2020-07-15 03:07:21}
\Level{B}
\SubArea{Logarithmic equations and inequalities}
\LANGen{
\Question{
Solve
\[ \log_2(x)-\log_4(x)=1\text{ .} \]}
{\( x=4 \)}{\( x=2 \)}{\( x_1=\frac12\text{, }x_2=4 \)}{\( x_1=-1\text{, }x_2=2 \)}{}{}}
\LANGpl{
\Question{
Rozwiąż
\[ \log_2(x)-\log_4(x)=1\text{ .} \]}
{\( x=4 \)}{\( x=2 \)}{\( x_1=\frac12\text{, }x_2=4 \)}{\( x_1=-1\text{, }x_2=2 \)}{}{}}
\LANGcs{
\Question{
Řešte rovnici
\[ \log_2(x)-\log_4(x)=1\text{ .} \]}
{\( x=4 \)}{\( x=2 \)}{\( x_1=\frac12\text{, }x_2=4 \)}{\( x_1=-1\text{, }x_2=2 \)}{}{}}
\LANGsk{
\Question{
Riešte rovnicu
\[ \log_2(x)-\log_4(x)=1\text{ .} \]}
{\( x=4 \)}{\( x=2 \)}{\( x_1=\frac12\text{, }x_2=4 \)}{\( x_1=-1\text{, }x_2=2 \)}{}{}}
\LANGes{
\Question{
Resuelve.
\[ \log_2(x)-\log_4(x)=1\text{ .} \]}
{\( x=4 \)}{\( x=2 \)}{\( x_1=\frac12\text{, }x_2=4 \)}{\( x_1=-1\text{, }x_2=2 \)}{}{}}
\End

\Data\ID{1003158802}
\Updated{2020-07-15 03:07:21}
\Level{B}
\SubArea{Logarithmic equations and inequalities}
\LANGen{
\Question{
Find the solution set of the following equation.
\[ \log_{0.2}(x)-4\log_{0.04}(x)=\log_5(x) \]}
{\( (0;\infty) \)}{\( \mathbb{R} \)}{\( \emptyset \)}{\( [ 0;\infty) \)}{}{}}
\LANGpl{
\Question{
Wyznacz zbiór rozwiązań podanego równania.
\[ \log_{0{,}2}(x)-4\log_{0{,}04}(x)=\log_5(x) \]}
{\( (0;\infty) \)}{\( \mathbb{R} \)}{\( \emptyset \)}{\( \langle 0;\infty) \)}{}{}}
\LANGcs{
\Question{
Najděte množinu řešení dané rovnice.
\[ \log_{0{,}2}(x)-4\log_{0{,}04}(x)=\log_5(x) \]}
{\( (0;\infty) \)}{\( \mathbb{R} \)}{\( \emptyset \)}{\( \langle 0;\infty) \)}{}{}}
\LANGsk{
\Question{
Nájdite množinu riešení danej rovnice.
\[ \log_{0{,}2}(x)-4\log_{0{,}04}(x)=\log_5(x) \]}
{\( (0;\infty) \)}{\( \mathbb{R} \)}{\( \emptyset \)}{\( \langle 0;\infty) \)}{}{}}
\LANGes{
\Question{
Halla el conjunto de soluciones de la siguiente ecuación. 
\[ \log_{0.2}(x)-4\log_{0.04}(x)=\log_5(x) \]}
{\( (0;\infty) \)}{\( \mathbb{R} \)}{\( \emptyset \)}{\( [ 0;\infty) \)}{}{}}
\End

\Data\ID{1003158803}
\Updated{2020-07-15 03:07:21}
\Level{B}
\SubArea{Logarithmic equations and inequalities}
\LANGen{
\Question{
Solve
\[ \log_{\frac12}(x)+2=3\log_2(x) \text{ .} \]}
{\( x=\sqrt2 \)}{\( x=2 \)}{\( x=-1 \)}{Equation has no solution.}{}{}}
\LANGpl{
\Question{
Rozwiąż
\[ \log_{\frac12}(x)+2=3\log_2(x) \text{ .} \]}
{\( x=\sqrt2 \)}{\( x=2 \)}{\( x=-1 \)}{Równanie nie ma rozwiązania.}{}{}}
\LANGcs{
\Question{
Řešte rovnici
\[ \log_{\frac12}(x)+2=3\log_2(x) \text{ .} \]}
{\( x=\sqrt2 \)}{\( x=2 \)}{\( x=-1 \)}{Rovnice nemá řešení.}{}{}}
\LANGsk{
\Question{
Riešte rovnicu
\[ \log_{\frac12}(x)+2=3\log_2(x) \text{ .} \]}
{\( x=\sqrt2 \)}{\( x=2 \)}{\( x=-1 \)}{Rovnica nemá riešenie.}{}{}}
\LANGes{
\Question{
Resuelve.
\[ \log_{\frac12}(x)+2=3\log_2(x) \text{ .} \]}
{\( x=\sqrt2 \)}{\( x=2 \)}{\( x=-1 \)}{La ecuación no tiene solución.}{}{}}
\End

\Data\ID{1003158804}
\Updated{2020-07-15 03:07:21}
\Level{B}
\SubArea{Logarithmic equations and inequalities}
\LANGen{
\Question{
Which of the following statements is not true about the equation? 
\[ \log_x16+\log_{\frac1x}4=2 \]}
{The solution is an odd number.}{The solution is \( x=2 \).}{The solution is an even number.}{The solution is a prime number.}{}{}}
\LANGpl{
\Question{
Które z poniższych stwierdzeń  nie jest prawdziwe w odniesieniu do podanego równania? 
\[ \log_x16+\log_{\frac1x}4=2 \]}
{Rozwiązaniem jest liczba nieparzysta.}{Rozwiązaniem jest  \( x=2 \).}{Rozwiązaniem jest liczba parzysta.}{Rozwiązaniem jest liczba pierwsza.}{}{}}
\LANGcs{
\Question{
Které z následujících tvrzení o rovnici  není  pravdivé?
\[ \log_x16+\log_{\frac1x}4=2 \]}
{Řešením rovnice je liché číslo.}{Řešením rovnice je \( x=2 \).}{Řešením rovnice je sudé číslo.}{Řešením rovnice je prvočíslo.}{}{}}
\LANGsk{
\Question{
Ktoré z následujúcich tvrdení o rovnici  nie je pravdivé?
\[ \log_x16+\log_{\frac1x}4=2 \]}
{Riešením rovnice je nepárne číslo.}{Riešením rovnice je \( x=2 \).}{Riešením rovnice je párne číslo.}{Riešením rovnice je prvočíslo.}{}{}}
\LANGes{
\Question{
¿Cuál de los siguientes enunciados de la ecuación no es verdadero?  
\[ \log_x16+\log_{\frac1x}4=2 \]}
{La solución es un número impar.}{La solución es \( x=2 \).}{La solución es un número par.}{La solución es un número primo.}{}{}}
\End

\Data\ID{1003158805}
\Updated{2020-07-15 03:07:21}
\Level{B}
\SubArea{Logarithmic equations and inequalities}
\LANGen{
\Question{
Which of the following statements about the equation is not true? 
\[ \log_{x^2}4+\log_x2+\log_{\frac1x}1=1 \]}
{The solution is a prime number.}{The solution is an even number.}{The solution is a positive number.}{The solution is an integer.}{}{}}
\LANGpl{
\Question{
Które z poniższych stwierdzeń nie jest prawdziwe w odniesieniu do podanego równania? 
\[ \log_{x^2}4+\log_x2+\log_{\frac1x}1=1 \]}
{Rozwiązaniem jest liczba pierwsza.}{Rozwiązaniem jest liczba parzysta.}{Rozwiązaniem jest liczba dodatnia.}{Rozwiązaniem jest liczba całkowita.}{}{}}
\LANGcs{
\Question{
Které z následujících tvrzení o rovnici  není  pravdivé?
\[ \log_{x^2}4+\log_x2+\log_{\frac1x}1=1 \]}
{Řešením rovnice je prvočíslo.}{Řešením rovnice je sudé číslo.}{Řešením rovnice je kladné číslo.}{Řešením rovnice je celé číslo.}{}{}}
\LANGsk{
\Question{
Ktoré z následujúcich tvrdení o rovnici  nie je pravdivé?
\[ \log_{x^2}4+\log_x2+\log_{\frac1x}1=1 \]}
{Riešením rovnice je prvočíslo.}{Riešením rovnice je párne číslo.}{Riešením rovnice je kladné číslo.}{Riešením rovnice je celé číslo.}{}{}}
\LANGes{
\Question{
¿Cuál de los siguientes enunciados de la ecuación no es verdadero? 
\[ \log_{x^2}4+\log_x2+\log_{\frac1x}1=1 \]}
{La solución en un número primo.}{La solución es un número par.}{La solución es un número positivo.}{La solución es un número entero.}{}{}}
\End

\Data\ID{1003159001}
\Updated{2020-07-15 03:07:21}
\Level{B}
\SubArea{Logarithmic equations and inequalities}
\LANGen{
\Question{
Solve.
\[ 5^{-x}=7 \]}
{\( x=-\log_57 \)}{\( x=-\log_75 \)}{\( x=\log_57 \)}{The equation has no solution.}{}{}}
\LANGpl{
\Question{
Rozwiąż.
\[ 5^{-x}=7 \]}
{\( x=-\log_57 \)}{\( x=-\log_75 \)}{\( x=\log_57 \)}{Równanie nie ma rozwiązania.}{}{}}
\LANGcs{
\Question{
Řešte danou rovnici
\[ 5^{-x}=7 \]}
{\( x=-\log_57 \)}{\( x=-\log_75 \)}{\( x=\log_57 \)}{Rovnice nemá řešení.}{}{}}
\LANGsk{
\Question{
Riešte danú rovnicu.
\[ 5^{-x}=7\]}
{\( x=-\log_57 \)}{\( x=-\log_75 \)}{\( x=\log_57 \)}{Rovnica nemá riešenie.}{}{}}
\LANGes{
\Question{
Resuelve.
\[ 5^{-x}=7 \]}
{\( x=-\log_57 \)}{\( x=-\log_75 \)}{\( x=\log_57 \)}{La ecuación no tiene solución.}{}{}}
\End

\Data\ID{1003159002}
\Updated{2020-07-15 03:07:21}
\Level{B}
\SubArea{Logarithmic equations and inequalities}
\LANGen{
\Question{
Solve.
\[ 3^{x-1}=4 \]}
{\( x=\log_312 \)}{\( x=\log_{12}3 \)}{\( x=\log_37 \)}{\( x=\log4 \)}{}{}}
\LANGpl{
\Question{
Rozwiąż.
\[ 3^{x-1}=4 \]}
{\( x=\log_312 \)}{\( x=\log_{12}3 \)}{\( x=\log_37 \)}{\( x=\log4 \)}{}{}}
\LANGcs{
\Question{
Řešte danou rovnici.
\[ 3^{x-1}=4 \]}
{\( x=\log_312 \)}{\( x=\log_{12}3 \)}{\( x=\log_37 \)}{\( x=\log4 \)}{}{}}
\LANGsk{
\Question{
Riešte danú rovnicu.
\[ 3^{x-1}=4 \]}
{\( x=\log_312 \)}{\( x=\log_{12}3 \)}{\( x=\log_37 \)}{\( x=\log4 \)}{}{}}
\LANGes{
\Question{
Resuelve.
\[ 3^{x-1}=4 \]}
{\( x=\log_312 \)}{\( x=\log_{12}3 \)}{\( x=\log_37 \)}{\( x=\log4 \)}{}{}}
\End

\Data\ID{1003159003}
\Updated{2020-07-15 03:07:21}
\Level{B}
\SubArea{Logarithmic equations and inequalities}
\LANGen{
\Question{
Solve.
\[ 2^x=3^{2-x}\]}
{\( x=\log_69 \)}{\( x=\log_96 \)}{\( x=\log_59 \)}{\( x=1 \)}{}{}}
\LANGpl{
\Question{
Rozwiąż.
\[ 2^x=3^{2-x}\]}
{\( x=\log_69 \)}{\( x=\log_96 \)}{\( x=\log_59 \)}{\( x=1 \)}{}{}}
\LANGcs{
\Question{
Řešte danou rovnici.
\[ 2^x=3^{2-x}\]}
{\( x=\log_69 \)}{\( x=\log_96 \)}{\( x=\log_59 \)}{\( x=1 \)}{}{}}
\LANGsk{
\Question{
Riešte danú rovnicu.
\[ 2^x=3^{2-x}\]}
{\( x=\log_69 \)}{\( x=\log_96 \)}{\( x=\log_59 \)}{\( x=1 \)}{}{}}
\LANGes{
\Question{
Resuelve.
\[ 2^x=3^{2-x}\]}
{\( x=\log_69 \)}{\( x=\log_96 \)}{\( x=\log_59 \)}{\( x=1 \)}{}{}}
\End

\Data\ID{1003159004}
\Updated{2020-11-12 04:11:41}
\Level{B}
\SubArea{Logarithmic equations and inequalities}
\LANGen{
\Question{
Which of the given numbers is the sum of all the solutions of the equations \( \text{(1)} \) and \( \text{(2)} \)?
\[ \begin{aligned} 10^{x-1}&=2   	&\text{(1)} \\
2^{1-x}&=5^x 	&\text{(2)} \end{aligned} \]}
{\( \log40 \)}{\( \log22 \)}{\( \log12+\log_72 \)}{\( \log12-\log_72  \)}{}{}}
\LANGpl{
\Question{
Która z poniższych liczb jest sumą wszystkich rozwiązań podanego równania \( \text{(1)} \) i \( \text{(2)} \)?
\[ \begin{aligned} 10^{x-1}&=2   	&\text{(1)} \\
2^{1-x}&=5^x 	&\text{(2)} \end{aligned} \]}
{\( \log40 \)}{\( \log22 \)}{\( \log12+\log_72 \)}{\( \log12-\log_72  \)}{}{}}
\LANGcs{
\Question{
Které z daných čísel je součtem všech řešení následujících rovnic \( \text{(1)} \) a \( \text{(2)} \)?
\[ \begin{aligned} 10^{x-1}&=2   	&\text{(1)} \\
2^{1-x}&=5^x 	&\text{(2)} \end{aligned} \]}
{\( \log40 \)}{\( \log22 \)}{\( \log12+\log_72 \)}{\( \log12-\log_72  \)}{}{}}
\LANGsk{
\Question{
Určte súčet všetkých koreňov rovníc \( \text{(1)} \) a \( \text{(2)} \).
\[ \begin{aligned} 10^{x-1}&=2   	&\text{(1)} \\
2^{1-x}&=5^x 	&\text{(2)} \end{aligned} \]}
{\( \log40 \)}{\( \log22 \)}{\( \log12+\log_72 \)}{\( \log12-\log_72  \)}{}{}}
\LANGes{
\Question{
Define la suma de todas las soluciones de las ecuaciones \( \text{(1)} \) y \( \text{(2)} \).
\[ \begin{aligned} 10^{x-1}&=2   	&\text{(1)} \\
2^{1-x}&=5^x 	&\text{(2)} \end{aligned} \]}
{\( \log40 \)}{\( \log22 \)}{\( \log12+\log_72 \)}{\( \log12-\log_72  \)}{}{}}
\End

\Data\ID{1003159005}
\Updated{2020-11-12 04:11:11}
\Level{B}
\SubArea{Logarithmic equations and inequalities}
\LANGen{
\Question{
Which of the given numbers is the product of all the solutions of the following exponential equation?
\[ x^{\log_3x} =x \]}
{\( 3 \)}{\( 0 \)}{\( 1 \)}{\( \frac13 \)}{}{}}
\LANGpl{
\Question{
Która z poniższych liczb jest iloczynem wszystkich rozwiązań podanego równania wykładniczego? 
\[ x^{\log_3x} =x \]}
{\( 3 \)}{\( 0 \)}{\( 1 \)}{\( \frac13 \)}{}{}}
\LANGcs{
\Question{
Které z daných čísel je součinem všech řešení následující rovnice?
\[ x^{\log_3x} =x \]}
{\( 3 \)}{\( 0 \)}{\( 1 \)}{\( \frac13 \)}{}{}}
\LANGsk{
\Question{
Určte súčin všetkých koreňov danej exponenciálnej rovnice.
\[ x^{\log_3x} =x \]}
{\( 3 \)}{\( 0 \)}{\( 1 \)}{\( \frac13 \)}{}{}}
\LANGes{
\Question{
Identifica el producto de todas las raíces de la siguiente ecuación exponencial. 
\[ x^{\log_3x} =x \]}
{\( 3 \)}{\( 0 \)}{\( 1 \)}{\( \frac13 \)}{}{}}
\End

\Data\ID{1003162701}
\Updated{2020-11-12 04:11:47}
\Level{B}
\SubArea{Logarithmic equations and inequalities}
\LANGen{
\Question{
Find the solution set of the following equation.
\[ \log_3x+\frac3{\log_3x}=4 \]}
{\( \{3;27\} \)}{\( \{1;3\} \)}{\( \{-3;-1\} \)}{\( \left\{\frac1{27};\frac13\right\} \)}{}{}}
\LANGpl{
\Question{
Wyznacz zbiór rozwiązań podanego równania. 
\[ \log_3x+\frac3{\log_3x}=4 \]}
{\( \{3;27\} \)}{\( \{1;3\} \)}{\( \{-3;-1\} \)}{\( \left\{\frac1{27};\frac13\right\} \)}{}{}}
\LANGcs{
\Question{
Určete množinu všech řešení následující rovnice.
\[ \log_3x+\frac3{\log_3x}=4 \]}
{\( \{3;27\} \)}{\( \{1;3\} \)}{\( \{-3;-1\} \)}{\( \left\{\frac1{27};\frac13\right\} \)}{}{}}
\LANGsk{
\Question{
Určte množinu koreňov danej rovnice.
\[ \log_3x+\frac3{\log_3x}=4 \]}
{\( \{3;27\} \)}{\( \{1;3\} \)}{\( \{-3;-1\} \)}{\( \left\{\frac1{27};\frac13\right\} \)}{}{}}
\LANGes{
\Question{
Halla el conjunto de soluciones de la siguiente ecuación. 
\[ \log_3x+\frac3{\log_3x}=4 \]}
{\( \{3;27\} \)}{\( \{1;3\} \)}{\( \{-3;-1\} \)}{\( \left\{\frac1{27};\frac13\right\} \)}{}{}}
\End

\Data\ID{1003162702}
\Updated{2020-11-12 04:11:20}
\Level{B}
\SubArea{Logarithmic equations and inequalities}
\LANGen{
\Question{
Which of the given numbers is the sum of all the solutions of the following equation?
\[ \log_2^2 x-3\log_2x+2=0 \]}
{\( 6 \)}{\( 3 \)}{\( -3 \)}{\( \frac34 \)}{}{}}
\LANGpl{
\Question{
Która z poniższych liczb jest sumą wszystkich rozwiązań podanego równania?
\[ \log_2^2 x-3\log_2x+2=0 \]}
{\( 6 \)}{\( 3 \)}{\( -3 \)}{\( \frac34 \)}{}{}}
\LANGcs{
\Question{
Které z daných čísel je součtem všech řešení následující rovnice?
\[ \log_2^2 x-3\log_2x+2=0. \]}
{\( 6 \)}{\( 3 \)}{\( -3 \)}{\( \frac34 \)}{}{}}
\LANGsk{
\Question{
Nájdite súčet všetkých koreňov danej rovnice
\[ \log_2^2 x-3\log_2x+2=0. \]}
{\( 6 \)}{\( 3 \)}{\( -3 \)}{\( \frac34 \)}{}{}}
\LANGes{
\Question{
Halla la suma de todas las raíces de la siguiente ecuación. 
\[ \log_2^2 x-3\log_2x+2=0 \]}
{\( 6 \)}{\( 3 \)}{\( -3 \)}{\( \frac34 \)}{}{}}
\End

\Data\ID{1003162703}
\Updated{2020-07-15 03:07:21}
\Level{B}
\SubArea{Logarithmic equations and inequalities}
\LANGen{
\Question{
How many solutions does the following equation have?
\[ \ln x^2=\ln^2 x-3 \]}
{exactly two positive solutions}{exactly two solutions, one  positive and one negative}{no solutions}{exactly one solution}{}{}}
\LANGpl{
\Question{
Ile rozwiązań ma podane równanie? 
\[ \ln x^2=\ln^2 x-3 \]}
{dokładnie dwa dodatnie rozwiązania}{dokładnie dwa rozwiązania - jedno dodatnie, drugie ujemne}{nie ma rozwiązania}{dokładnie jedno rozwiązanie}{}{}}
\LANGcs{
\Question{
Kolik řešení má daná rovnice?
\[ \ln x^2=\ln^2 x-3 \]}
{právě dvě kladné řešení}{právě dvě řešení, jedno kladné a jedno záporné}{nemá řešení}{právě jedno řešení}{}{}}
\LANGsk{
\Question{
Koľko riešení má daná rovnica?
\[ \ln x^2=\ln^2 x-3 \]}
{práve dve kladné riešenia}{práve dve riešenia, jedno kladné a jedno záporné}{nemá riešenie}{práve jedno riešenie}{}{}}
\LANGes{
\Question{
¿Cuántas soluciones tiene la siguiente ecuación? 
\[ \ln x^2=\ln^2 x-3 \]}
{Exactamente dos soluciones positivas.}{Exactamente dos soluciones, una positiva y una negativa.}{Sin solución.}{Exactamente una solución.}{}{}}
\End

\Data\ID{1003162704}
\Updated{2020-07-15 03:07:21}
\Level{B}
\SubArea{Logarithmic equations and inequalities}
\LANGen{
\Question{
Which of the following statements about the given equation is true? 
\[ \log_4(x-1)^2=3-\frac1{\log_4(x-1)} \]}
{The solution set consists of two prime numbers.}{The solution set is \( \left\{\frac12;1\right\} \).}{The solution set is an empty set.}{The equation has exactly one solution.}{None of the above statements is true.}{}}
\LANGpl{
\Question{
Które z poniższych stwierdzeń o podanym równaniu jest prawdziwe? 
\[ \log_4(x-1)^2=3-\frac1{\log_4(x-1)} \]}
{Zbiorem rozwiązań są dwie liczby pierwsze.}{Zbiorem rozwiązań jest  \( \left\{\frac12;1\right\} \).}{Zbiorem rozwiązań jest zbiór pusty.}{Równanie ma dokładnie jedno rozwiązanie.}{Żadne z powyższych stwierdzeń nie jest prawdziwe.}{}}
\LANGcs{
\Question{
Které z následujících tvrzení o dané rovnici je pravdivé?
\[ \log_4(x-1)^2=3-\frac1{\log_4(x-1)} \]}
{Řešením jsou právě dvě prvočísla.}{Množina řešení je \( \left\{\frac12;1\right\} \).}{Řešením je prázdná množina.}{Rovnice má právě jedno řešení.}{Žádné tvrzení není  pravdivé.}{}}
\LANGsk{
\Question{
Ktoré z nasledujúcich tvrdení o danej rovnici je pravdivé?
\[ \log_4(x-1)^2=3-\frac1{\log_4(x-1)} \]}
{Riešením sú práve dve prvočísla.}{Množina riešení je \( \left\{\frac12;1\right\} \).}{Riešením je prázdna množina.}{Rovnica má práve jedno riešenie.}{Žiadne tvrdenie nie je pravdivé.}{}}
\LANGes{
\Question{
¿Cuál de los siguientes enunciados de la siguiente ecuación es verdadero?  
\[ \log_4(x-1)^2=3-\frac1{\log_4(x-1)} \]}
{El conjunto de soluciones son dos números primos.}{El conjunto de soluciones es \( \left\{\frac12;1\right\} \).}{El conjunto de soluciones es un conjunto vacío.}{La ecuación tiene única solución.}{Ninguno de los enunciados es verdadero.}{}}
\End

\Data\ID{1003177802}
\Updated{2020-07-15 03:07:21}
\Level{B}
\SubArea{Logarithmic equations and inequalities}
\LANGen{
\Question{
Choose the domain of the expression \( \ln\!\left(-|3-2x|+6\right) \).}
{\( \left(-\frac32;\frac92\right) \)}{\( \left(-\infty;-\frac32\right)\cup\left(\frac92;\infty\right) \)}{\( \left(-\infty;\frac32\right)\cup\left(\frac92;\infty\right) \)}{\( \left(\frac92;\infty\right) \)}{}{}}
\LANGpl{
\Question{
Wybierz dziedzinę wyrażenia \( \ln\!\left(-|3-2x|+6\right) \).}
{\( \left(-\frac32;\frac92\right) \)}{\( \left(-\infty;-\frac32\right)\cup\left(\frac92;\infty\right) \)}{\( \left(-\infty;\frac32\right)\cup\left(\frac92;\infty\right) \)}{\( \left(\frac92;\infty\right) \)}{}{}}
\LANGcs{
\Question{
Určete definiční obor výrazu \( \ln\!\left(-|3-2x|+6\right) \).}
{\( \left(-\frac32;\frac92\right) \)}{\( \left(-\infty;-\frac32\right)\cup\left(\frac92;\infty\right) \)}{\( \left(-\infty;\frac32\right)\cup\left(\frac92;\infty\right) \)}{\( \left(\frac92;\infty\right) \)}{}{}}
\LANGsk{
\Question{
Určte definičný obor výrazu \( \ln\!\left(-|3-2x|+6\right) \).}
{\( \left(-\frac32;\frac92\right) \)}{\( \left(-\infty;-\frac32\right)\cup\left(\frac92;\infty\right) \)}{\( \left(-\infty;\frac32\right)\cup\left(\frac92;\infty\right) \)}{\( \left(\frac92;\infty\right) \)}{}{}}
\LANGes{
\Question{
Elige el dominio de la expresión \( \ln\!\left(-|3-2x|+6\right) \).}
{\( \left(-\frac32;\frac92\right) \)}{\( \left(-\infty;-\frac32\right)\cup\left(\frac92;\infty\right) \)}{\( \left(-\infty;\frac32\right)\cup\left(\frac92;\infty\right) \)}{\( \left(\frac92;\infty\right) \)}{}{}}
\End

\Data\ID{9000003805}
\Updated{2020-07-15 03:07:34}
\Level{B}
\SubArea{Logarithmic equations and inequalities}
\LANGen{
\Question{
Solve the following equation. 
\[ 
 \log x^{2}\cdot \log \sqrt{x} -\log \frac{1} 
{x} = 2
\]}
{\(x_{1} = \frac{1} 
{100}\),
\(x_{2} = 10\)}{\(x_{1} = -2\),
\(x_{2} = 1\)}{\(x_{1} = - \frac{1}
{100}\),
\(x_{2} = 10\)}{\(x_{1} = -1\),
\(x_{2} = 2\)}{}{}}
\LANGpl{
\Question{
Rozwiąż następujące równanie.
\[ 
 \log x^{2}\cdot \log \sqrt{x} -\log \frac{1} 
{x} = 2
\]}
{\(x_{1} = \frac{1} 
{100}\),
\(x_{2} = 10\)}{\(x_{1} = -2\),
\(x_{2} = 1\)}{\(x_{1} = - \frac{1}
{100}\),
\(x_{2} = 10\)}{\(x_{1} = -1\),
\(x_{2} = 2\)}{}{}}
\LANGcs{
\Question{
Řešením logaritmické rovnice \(\log x^{2}\cdot \log \sqrt{x} -\log \frac{1} 
{x} = 2\) 
s neznámou \(x\in \mathbb{R}\) 
je:}
{\(x_{1} = \frac{1} 
{100}\),
\(x_{2} = 10\)}{\(x_{1} = -2\),
\(x_{2} = 1\)}{\(x_{1} = - \frac{1}
{100}\),
\(x_{2} = 10\)}{\(x_{1} = -1\),
\(x_{2} = 2\)}{}{}}
\LANGsk{
\Question{
Riešením rovnice  
\(
 \log x^{2}\cdot \log \sqrt{x} -\log \frac{1} 
{x} = 2
\)  sú korene:}
{\(x_{1} = \frac{1} 
{100}\),
\(x_{2} = 10\)}{\(x_{1} = -2\),
\(x_{2} = 1\)}{\(x_{1} = - \frac{1}
{100}\),
\(x_{2} = 10\)}{\(x_{1} = -1\),
\(x_{2} = 2\)}{}{}}
\LANGes{
\Question{
Resuelve la siguiente ecuación.
\[ 
 \log x^{2}\cdot \log \sqrt{x} -\log \frac{1} 
{x} = 2
\]}
{\(x_{1} = \frac{1} 
{100}\),
\(x_{2} = 10\)}{\(x_{1} = -2\),
\(x_{2} = 1\)}{\(x_{1} = - \frac{1}
{100}\),
\(x_{2} = 10\)}{\(x_{1} = -1\),
\(x_{2} = 2\)}{}{}}
\End

\Data\ID{9000003806}
\Updated{2020-07-15 03:07:34}
\Level{B}
\SubArea{Logarithmic equations and inequalities}
\LANGen{
\Question{
In the following list identify an equation such that neither
\(x = 5\) nor
\(x = 3\) is the
solution of this equation.}
{\(\log _{3}(1 - x) =\log _{3}(x + 16 - x^{2})\)}{\(\log (54 - x^{3}) = 3\cdot \log x\)}{\(\log _{5}(x^{2} - 17) =\log _{5}(x + 3)\)}{\(\log (x - 2) -\log (4 - x) = 1 -\log (13 - x)\)}{}{}}
\LANGpl{
\Question{
Dla którego z poniższych równań rozwiązaniem nie jest ani 
\(x = 5\) ani
\(x = 3\)?}
{\(\log _{3}(1 - x) =\log _{3}(x + 16 - x^{2})\)}{\(\log (54 - x^{3}) = 3\cdot \log x\)}{\(\log _{5}(x^{2} - 17) =\log _{5}(x + 3)\)}{\(\log (x - 2) -\log (4 - x) = 1 -\log (13 - x)\)}{}{}}
\LANGcs{
\Question{
Určete, která z daných logaritmických rovnic nemá řešení ani
\(x = 5\) ani
\(x = 3\).}
{\(\log _{3}(1 - x) =\log _{3}(x + 16 - x^{2})\)}{\(\log (54 - x^{3}) = 3\cdot \log x\)}{\(\log _{5}(x^{2} - 17) =\log _{5}(x + 3)\)}{\(\log (x - 2) -\log (4 - x) = 1 -\log (13 - x)\)}{}{}}
\LANGsk{
\Question{
Vyberte rovnicu, ktorej riešením nie je koreň
\(x = 5\) ani koreň
\(x = 3\).}
{\(\log _{3}(1 - x) =\log _{3}(x + 16 - x^{2})\)}{\(\log (54 - x^{3}) = 3\cdot \log x\)}{\(\log _{5}(x^{2} - 17) =\log _{5}(x + 3)\)}{\(\log (x - 2) -\log (4 - x) = 1 -\log (13 - x)\)}{}{}}
\LANGes{
\Question{
Identifica cuál de las siguientes ecuaciones no tiene solución ni 
\(x = 5\) ni
\(x = 3\)}
{\(\log _{3}(1 - x) =\log _{3}(x + 16 - x^{2})\)}{\(\log (54 - x^{3}) = 3\cdot \log x\)}{\(\log _{5}(x^{2} - 17) =\log _{5}(x + 3)\)}{\(\log (x - 2) -\log (4 - x) = 1 -\log (13 - x)\)}{}{}}
\End

\Data\ID{9000003808}
\Updated{2020-07-15 03:07:34}
\Level{B}
\SubArea{Logarithmic equations and inequalities}
\LANGen{
\Question{
Identify one statement which is true for the following equation. 
\[ 
 \log (x - 13) -\log (x - 3) = 1 -\log 2
\]}
{The equation does not have a solution.}{The equation has two solutions.}{The equation has a unique solution. This solution is a noninteger rational
number.}{The solution is \(x = 0\).}{The equation has a unique solution. This solution is a positive integer.}{The equation has a unique solution. This solution is a negative integer.}}
\LANGpl{
\Question{
Jedno z poniższych zdań dotyczących podanego równania jest prawdziwe. Które?
\[ 
 \log (x - 13) -\log (x - 3) = 1 -\log 2
\]}
{Równanie nie ma rozwiązania.}{Równanie ma dwa rozwiązania.}{Równanie ma tylko jedno rozwiązanie. Tym rozwiązaniem jest liczba wymierna niecałkowita.}{Rozwiązaniem jest \(x = 0\).}{Równanie ma tylko jedno rozwiązanie. Tym rozwiązaniem jest liczba całkowita dodatnia.}{Równanie ma tylko jedno rozwiązanie. Tym rozwiązaniem jest liczba całkowita ujemna.}}
\LANGcs{
\Question{
Uvažujme rovnici
\[\log (x - 13) -\log (x - 3) = 1 -\log 2\] s
neznámou \(x\in \mathbb{R}\).
Vyberte, které z následujících tvrzení je pravdivé.}
{Rovnice nemá řešení.}{Rovnice má právě dvě řešení.}{Rovnice má právě jedno řešení. Toto řešení je racionálním číslem a není
celým číslem.}{Rovnice má řešení \(x=0\).}{Rovnice má právě jedno řešení. Toto řešení je přirozeným číslem.}{Rovnice má právě jedno řešení. Toto řešení je záporným celým číslem.}}
\LANGsk{
\Question{
Je daná rovnica
\[\log (x - 13) -\log (x - 3) = 1 -\log 2\] s
neznámou \(x\in \mathbb{R}\).
Vyberte, ktoré z následujúcich tvrdení o rovnici je pravdivé.}
{Rovnica nemá riešenie.}{Rovnica má práve dve riešenia.}{Rovnica má práva jedno riešenie. Toto riešenie je racionálne číslo a zároveň nie je
celé číslo.}{Riešením rovnice je koreň \(x=0\).}{Rovnica má práve jedno riešenie. Toto riešenie je prirodzené číslo.}{Rovnica má práve jedno riešenie. Toto riešenie je záporné celé číslo.}}
\LANGes{
\Question{
Identifica cuál de los enunciados de la siguiente ecuación es verdadero.  
\[ 
 \log (x - 13) -\log (x - 3) = 1 -\log 2
\]}
{La ecuación no tiene solución.}{La ecuación tiene exactamente dos soluciones.}{La ecuación tiene única solución. La solución es un número racional y no es un número entero.}{La solución es \(x = 0\).}{La ecuación tiene una única solución. La solución es un número natural positivo.}{La ecuación tiene una única solución. La solución es un número entero negativo.}}
\End

\Data\ID{1003162801}
\Updated{2020-07-15 03:07:21}
\Level{C}
\SubArea{Logarithmic equations and inequalities}
\LANGen{
\Question{
Find the solution set of the following inequality.
\[ \log_{0.3}(x)\geq \log_{0.3}3 \]}
{\( (0;3] \)}{\( (-\infty;3] \)}{\( [3;\infty) \)}{\( [0;3] \)}{}{}}
\LANGpl{
\Question{
Wyznacz zbiór rozwiązań podanej nierówności.
\[ \log_{0{,}3}(x)\geq \log_{0{,}3}3 \]}
{\( (0;3\rangle \)}{\( (-\infty;3\rangle \)}{\( \langle3;\infty) \)}{\( \langle0;3\rangle \)}{}{}}
\LANGcs{
\Question{
Vyberte množinu, která je řešením dané nerovnice.
\[ \log_{0{,}3}(x)\geq \log_{0{,}3}3 \]}
{\( (0;3\rangle \)}{\( (-\infty;3\rangle \)}{\( \langle3;\infty) \)}{\( \langle0;3\rangle \)}{}{}}
\LANGsk{
\Question{
Vyberte množinu, ktorá je riešením danej nerovnice.
\[ \log_{0{,}3}(x)\geq \log_{0{,}3}3 \]}
{\( (0;3\rangle \)}{\( (-\infty;3\rangle \)}{\( \langle3;\infty) \)}{\( \langle0;3\rangle \)}{}{}}
\LANGes{
\Question{
Halla el conjunto de soluciones de la siguiente inecuación. 
\[ \log_{0.3}(x)\geq \log_{0.3}3 \]}
{\( (0;3] \)}{\( (-\infty;3] \)}{\( [3;\infty) \)}{\( [0;3] \)}{}{}}
\End

\Data\ID{1003162802}
\Updated{2020-07-15 03:07:21}
\Level{C}
\SubArea{Logarithmic equations and inequalities}
\LANGen{
\Question{
Find the solution set of the following inequality.
\[ \log_{\frac12}(x) \leq 1 \]}
{\( \left[\frac12;\infty\right) \)}{\( \left(-\infty;\frac12\right] \)}{\( \left(0;\frac12\right] \)}{\( [1;\infty) \)}{}{}}
\LANGpl{
\Question{
Wyznacz zbiór rozwiązań podanej nierówności. 
\[ \log_{\frac12}(x) \leq 1 \]}
{\( \left\langle\frac12;\infty\right) \)}{\( \left(-\infty;\frac12\right\rangle \)}{\( \left(0;\frac12\right\rangle \)}{\( \langle1;\infty) \)}{}{}}
\LANGcs{
\Question{
Vyberte množinu, která je řešením dané nerovnice.
\[ \log_{\frac12}(x) \leq 1 \]}
{\( \left\langle\frac12;\infty\right) \)}{\( \left(-\infty;\frac12\right\rangle \)}{\( \left(0;\frac12\right\rangle \)}{\( \langle1;\infty) \)}{}{}}
\LANGsk{
\Question{
Vyberte množinu, ktorá je riešením danej nerovnice.
\[ \log_{\frac12}(x) \leq 1 \]}
{\( \left\langle\frac12;\infty\right) \)}{\( \left(-\infty;\frac12\right\rangle \)}{\( \left(0;\frac12\right\rangle \)}{\( \langle1;\infty) \)}{}{}}
\LANGes{
\Question{
Halla el conjunto de soluciones de la siguiente inecuación. 
\[ \log_{\frac12}(x) \leq 1 \]}
{\( \left[\frac12;\infty\right) \)}{\( \left(-\infty;\frac12\right] \)}{\( \left(0;\frac12\right] \)}{\( [1;\infty) \)}{}{}}
\End

\Data\ID{1003162803}
\Updated{2020-07-15 03:07:21}
\Level{C}
\SubArea{Logarithmic equations and inequalities}
\LANGen{
\Question{
Find the solution set of the following inequality.
\[ \log_4(2x-1)\leq 0 \]}
{\( \left(\frac12;1\right] \)}{\( \left(-\infty;1\right] \)}{\( \left\{\frac12\right\} \)}{\( \left(-\infty;\frac12\right] \)}{}{}}
\LANGpl{
\Question{
Wyznacz zbiór rozwiązań podanej nierówności. 
\[ \log_4(2x-1)\leq 0 \]}
{\( \left(\frac12;1\right\rangle \)}{\( \left(-\infty;1\right\rangle \)}{\( \left\{\frac12\right\} \)}{\( \left(-\infty;\frac12\right\rangle \)}{}{}}
\LANGcs{
\Question{
Vyberte množinu, která je řešením dané nerovnice.
\[ \log_4(2x-1)\leq 0 \]}
{\( \left(\frac12;1\right\rangle \)}{\( \left(-\infty;1\right\rangle \)}{\( \left\{\frac12\right\} \)}{\( \left(-\infty;\frac12\right\rangle \)}{}{}}
\LANGsk{
\Question{
Vyberte množinu, ktorá je riešením danej nerovnice.
\[ \log_4(2x-1)\leq 0 \]}
{\( \left(\frac12;1\right\rangle \)}{\( \left(-\infty;1\right\rangle \)}{\( \left\{\frac12\right\} \)}{\( \left(-\infty;\frac12\right\rangle \)}{}{}}
\LANGes{
\Question{
Halla el conjunto de soluciones de la siguiente inecuación. 
\[ \log_4(2x-1)\leq 0 \]}
{\( \left(\frac12;1\right] \)}{\( \left(-\infty;1\right] \)}{\( \left\{\frac12\right\} \)}{\( \left(-\infty;\frac12\right] \)}{}{}}
\End

\Data\ID{1003162804}
\Updated{2020-07-15 03:07:21}
\Level{C}
\SubArea{Logarithmic equations and inequalities}
\LANGen{
\Question{
Find the solution set of the following inequality.
\[ \log_{16} x^2 < \frac12 \]}
{\( (-2;0)\cup(0;2) \)}{\( (-\infty;-2)\cup(2;\infty) \)}{\( (0;2) \)}{\( (2;\infty) \)}{}{}}
\LANGpl{
\Question{
Wyznacz zbiór rozwiązań podanej nierówności.
\[ \log_{16} x^2 < \frac12 \]}
{\( (-2;0)\cup(0;2) \)}{\( (-\infty;-2)\cup(2;\infty) \)}{\( (0;2) \)}{\( (2;\infty) \)}{}{}}
\LANGcs{
\Question{
Vyberte množinu, která je řešením dané nerovnice.
\[ \log_{16} x^2 < \frac12 \]}
{\( (-2;0)\cup(0;2) \)}{\( (-\infty;-2)\cup(2;\infty) \)}{\( (0;2) \)}{\( (2;\infty) \)}{}{}}
\LANGsk{
\Question{
Vyberte množinu, ktorá je riešením danej nerovnice.
\[ \log_{16} x^2 < \frac12 \]}
{\( (-2;0)\cup(0;2) \)}{\( (-\infty;-2)\cup(2;\infty) \)}{\( (0;2) \)}{\( (2;\infty) \)}{}{}}
\LANGes{
\Question{
Halla el conjunto de soluciones de la siguiente inecuación. 
\[ \log_{16} x^2 < \frac12 \]}
{\( (-2;0)\cup(0;2) \)}{\( (-\infty;-2)\cup(2;\infty) \)}{\( (0;2) \)}{\( (2;\infty) \)}{}{}}
\End

\Data\ID{1003170901}
\Updated{2020-07-15 03:07:21}
\Level{C}
\SubArea{Logarithmic equations and inequalities}
\LANGen{
\Question{
Find the solution set of the following inequality.
\[ \log_3\frac{2x-2}{x+3}\leq 0 \]}
{\( (1;5] \)}{\( (-3;5] \)}{\( [-3;5] \)}{\( [-3;1] \)}{}{}}
\LANGpl{
\Question{
Wyznacz zbiór rozwiązań podanej nierówności. 
\[ \log_3\frac{2x-2}{x+3}\leq 0 \]}
{\( (1;5\rangle \)}{\( (-3;5\rangle \)}{\( \langle-3;5\rangle \)}{\( \langle-3;1\rangle \)}{}{}}
\LANGcs{
\Question{
Najděte množinu řešení dané nerovnice.
\[ \log_3\frac{2x-2}{x+3}\leq 0 \]}
{\( (1;5\rangle \)}{\( (-3;5\rangle \)}{\( \langle-3;5\rangle \)}{\( \langle-3;1\rangle \)}{}{}}
\LANGsk{
\Question{
Nájdite množinu riešení danej nerovnice.
\[ \log_3\frac{2x-2}{x+3}\leq 0 \]}
{\( (1;5\rangle \)}{\( (-3;5\rangle \)}{\( \langle-3;5\rangle \)}{\( \langle-3;1\rangle \)}{}{}}
\LANGes{
\Question{
Halla el conjunto de soluciones de la siguiente inecuación. 
\[ \log_3\frac{2x-2}{x+3}\leq 0 \]}
{\( (1;5] \)}{\( (-3;5] \)}{\( [-3;5] \)}{\( [-3;1] \)}{}{}}
\End

\Data\ID{1003170902}
\Updated{2020-07-15 03:07:21}
\Level{C}
\SubArea{Logarithmic equations and inequalities}
\LANGen{
\Question{
Find the solution set of the following inequality.
\[ \log_{0.1}\frac{2-x}{x+1}>0 \]}
{\( \left(\frac12;2\right) \)}{\( (-\infty;-1)\cup\left(\frac12;\infty\right) \)}{\( \left(-1;\frac12 \right) \)}{\( (-1;2) \)}{}{}}
\LANGpl{
\Question{
Wyznacz zbiór rozwiązań podanej nierówności. 
\[ \log_{0{,}1}\frac{2-x}{x+1}>0 \]}
{\( \left(\frac12;2\right) \)}{\( (-\infty;-1)\cup\left(\frac12;\infty\right) \)}{\( \left(-1;\frac12 \right) \)}{\( (-1;2) \)}{}{}}
\LANGcs{
\Question{
Najděte množinu řešení dané nerovnice.
\[ \log_{0{,}1}\frac{2-x}{x+1}>0 \]}
{\( \left(\frac12;2\right) \)}{\( (-\infty;-1)\cup\left(\frac12;\infty\right) \)}{\( \left(-1;\frac12 \right) \)}{\( (-1;2) \)}{}{}}
\LANGsk{
\Question{
Nájdite množinu riešení danej nerovnice.
\[ \log_{0{,}1}\frac{2-x}{x+1}>0 \]}
{\( \left(\frac12;2\right) \)}{\( (-\infty;-1)\cup\left(\frac12;\infty\right) \)}{\( \left(-1;\frac12 \right) \)}{\( (-1;2) \)}{}{}}
\LANGes{
\Question{
Halla el conjunto de soluciones de la siguiente inecuación. 
\[ \log_{0.1}\frac{2-x}{x+1}>0 \]}
{\( \left(\frac12;2\right) \)}{\( (-\infty;-1)\cup\left(\frac12;\infty\right) \)}{\( \left(-1;\frac12 \right) \)}{\( (-1;2) \)}{}{}}
\End

\Data\ID{1003170903}
\Updated{2020-07-15 03:07:21}
\Level{C}
\SubArea{Logarithmic equations and inequalities}
\LANGen{
\Question{
Find the solution set of the following inequality.
\[ \log_2^2(x-2)< \log_2(x-2) \]}
{\( (3;4) \)}{\( (0;1) \)}{\( (-\infty;4) \)}{\( (-\infty;1) \)}{}{}}
\LANGpl{
\Question{
Wyznacz zbiór rozwiązań podanej nierówności.
\[ \log_2^2(x-2)< \log_2(x-2) \]}
{\( (3;4) \)}{\( (0;1) \)}{\( (-\infty;4) \)}{\( (-\infty;1) \)}{}{}}
\LANGcs{
\Question{
Najdi množinu řešení dané nerovnice.
\[ \log_2^2(x-2)< \log_2(x-2) \]}
{\( (3;4) \)}{\( (0;1) \)}{\( (-\infty;4) \)}{\( (-\infty;1) \)}{}{}}
\LANGsk{
\Question{
Nájdite množinu riešení danej nerovnice.
\[ \log_2^2(x-2)< \log_2(x-2) \]}
{\( (3;4) \)}{\( (0;1) \)}{\( (-\infty;4) \)}{\( (-\infty;1) \)}{}{}}
\LANGes{
\Question{
Halla el conjunto de soluciones de la siguiente inecuación. 
\[ \log_2^2(x-2)< \log_2(x-2) \]}
{\( (3;4) \)}{\( (0;1) \)}{\( (-\infty;4) \)}{\( (-\infty;1) \)}{}{}}
\End

\Data\ID{1003170904}
\Updated{2020-07-15 03:07:21}
\Level{C}
\SubArea{Logarithmic equations and inequalities}
\LANGen{
\Question{
Find the solution set of the following inequality.
\[ \frac{\log x-1}{2+\log x}\geq1 \]}
{\( \left(0;\frac1{100}\right) \)}{\( \left(0;\frac1{100}\right] \)}{\( \left(-\infty;\frac1{100}\right) \)}{\( \left(-\infty;\frac1{100}\right] \)}{}{}}
\LANGpl{
\Question{
Wyznacz zbiór rozwiązań podanej nierówności. 
\[ \frac{\log x-1}{2+\log x}\geq1 \]}
{\( \left(0;\frac1{100}\right) \)}{\( \left(0;\frac1{100}\right\rangle \)}{\( \left(-\infty;\frac1{100}\right) \)}{\( \left(-\infty;\frac1{100}\right\rangle \)}{}{}}
\LANGcs{
\Question{
Najdi množinu řešení dané nerovnice.
\[ \frac{\log x-1}{2+\log x}\geq1 \]}
{\( \left(0;\frac1{100}\right) \)}{\( \left(0;\frac1{100}\right\rangle \)}{\( \left(-\infty;\frac1{100}\right) \)}{\( \left(-\infty;\frac1{100}\right\rangle \)}{}{}}
\LANGsk{
\Question{
Nájdite množinu riešení danej nerovnice.
\[ \frac{\log x-1}{2+\log x}\geq1 \]}
{\( \left(0;\frac1{100}\right) \)}{\( \left(0;\frac1{100}\right\rangle \)}{\( \left(-\infty;\frac1{100}\right) \)}{\( \left(-\infty;\frac1{100}\right\rangle \)}{}{}}
\LANGes{
\Question{
Halla el conjunto de soluciones de la siguiente inecuación. 
\[ \frac{\log x-1}{2+\log x}\geq1 \]}
{\( \left(0;\frac1{100}\right) \)}{\( \left(0;\frac1{100}\right] \)}{\( \left(-\infty;\frac1{100}\right) \)}{\( \left(-\infty;\frac1{100}\right] \)}{}{}}
\End

\Data\ID{1103162805}
\Updated{2020-07-15 03:07:10}
\Level{C}
\SubArea{Logarithmic equations and inequalities}
\LANGen{
\Question{
Choose the diagram which shows the number line with the correct solution set of the given inequality depicted in blue.
\[ 2\log_{0.1}(x-1) > \log_{0.1}4 \]}
{}{}{}{}{}{}}
\LANGpl{
\Question{
Wybierz wykres, który przedstawia zbiór rozwiązań podanej nierówności oznaczonej kolorem niebieskim. 
\[ 2\log_{0{,}1}(x-1) > \log_{0{,}1}4 \]}
{}{}{}{}{}{}}
\LANGcs{
\Question{
Rozhodněte, který z následujících intervalů znázorněných na číselné ose modrou barvou je řešením dané nerovnice.
\[ 2\log_{0{,}1}(x-1) > \log_{0{,}1}4 \]}
{}{}{}{}{}{}}
\LANGsk{
\Question{
Rozhodnite, ktorý z nasledujúcich intervalov znázornených na číselnej osi modrou farbou je riešením danej nerovnice.
\[ 2\log_{0{,}1}(x-1) > \log_{0{,}1}4 \]}
{}{}{}{}{}{}}
\LANGes{
\Question{
NA}
{}{}{}{}{}{}}

\IMGanswer{1103162805a.pdf}
\IMGanswer{1103162805b.pdf}
\IMGanswer{1103162805c.pdf}
\IMGanswer{1103162805d.pdf}\End

\Data\ID{1103162806}
\Updated{2020-07-15 03:07:10}
\Level{C}
\SubArea{Logarithmic equations and inequalities}
\LANGen{
\Question{
Choose the diagram which shows the number line with the correct solution of the given inequality depicted in blue.
\[ \log_{\frac13}(2-x) < \log_{\frac13}x+\log_{\frac13}3 \]}
{}{}{}{}{}{}}
\LANGpl{
\Question{
Wybierz wykres, który przedstawia zbiór rozwiązań podanej nierówności w kolorze niebieskim. 
\[ \log_{\frac13}(2-x) < \log_{\frac13}x+\log_{\frac13}3 \]}
{}{}{}{}{}{}}
\LANGcs{
\Question{
Rozhodněte, který z následujících intervalů znázorněných na číselné ose modrou barvou je řešením dané nerovnice.
\[ \log_{\frac13}(2-x) < \log_{\frac13}x+\log_{\frac13}3 \]}
{}{}{}{}{}{}}
\LANGsk{
\Question{
Rozhodnite, ktorý z nasledujúcich intervalov znázornených na číselnej osi modrou farbou je riešením danej nerovnice.
\[ \log_{\frac13}(2-x) < \log_{\frac13}x+\log_{\frac13}3 \]}
{}{}{}{}{}{}}
\LANGes{
\Question{
NA}
{}{}{}{}{}{}}

\IMGanswer{1103162806a.pdf}
\IMGanswer{1103162806b.pdf}
\IMGanswer{1103162806c.pdf}
\IMGanswer{1103162806d.pdf}\End

\Data\ID{9000003809}
\Updated{2020-10-13 02:10:31}
\Level{C}
\SubArea{Logarithmic equations and inequalities}
\LANGen{
\Question{
Find the solution set of the following inequality. 
\[ 
 \log _{0.5}(x^{2} - 2x) >\log _{
0.5}3
\]}
{\((-1;0)\cup (2;3)\)}{\((-\infty ;0)\cup (2;\infty )\)}{\((0;2)\)}{\((-\infty ;-1)\cup (0;2)\cup (3;\infty )\)}{\((-\infty ;-1)\cup (3;\infty )\)}{\((-1;3)\)}}
\LANGpl{
\Question{
Wyznacz zbiór rozwiązań podanej nierówności.
\[ 
 \log _{0.5}(x^{2} - 2x) >\log _{
0.5}3
\]}
{\((-1;0)\cup (2;3)\)}{\((-\infty ;0)\cup (2;\infty )\)}{\((0;2)\)}{\((-\infty ;-1)\cup (0;2)\cup (3;\infty )\)}{\((-\infty ;-1)\cup (3;\infty )\)}{\((-1;3)\)}}
\LANGcs{
\Question{
Najděte množinu řešení dané nerovnice.
\[\log _{0{,}5}(x^{2} - 2x) >\log _{0{,}5}3\]}
{\((-1;0)\cup (2;3)\)}{\((-\infty ;0)\cup (2;\infty )\)}{\((0;2)\)}{\((-\infty ;-1)\cup (0;2)\cup (3;\infty )\)}{\((-\infty ;-1)\cup (3;\infty )\)}{\((-1;3)\)}}
\LANGsk{
\Question{
Nájdite množinu riešení danej nerovnice.
\[\log _{0{,}5}(x^{2} - 2x) >\log _{0{,}5}3\]}
{\((-1;0)\cup (2;3)\)}{\((-\infty ;0)\cup (2;\infty )\)}{\((0;2)\)}{\((-\infty ;-1)\cup (0;2)\cup (3;\infty )\)}{\((-\infty ;-1)\cup (3;\infty )\)}{\((-1;3)\)}}
\LANGes{
\Question{
Halla el conjunto de soluciones de la siguiente inecuación.  
\[ 
 \log _{0.5}(x^{2} - 2x) >\log _{
0.5}3
\]}
{\((-1;0)\cup (2;3)\)}{\((-\infty ;0)\cup (2;\infty )\)}{\((0;2)\)}{\((-\infty ;-1)\cup (0;2)\cup (3;\infty )\)}{\((-\infty ;-1)\cup (3;\infty )\)}{\((-1;3)\)}}
\End

\Data\ID{9000004901}
\Updated{2020-07-15 03:07:34}
\Level{C}
\SubArea{Logarithmic equations and inequalities}
\LANGen{
\Question{
Solve the following inequality. 
\[ 
 \log _{0.3}x\geq \log _{0.3}5
\]}
{\(x\in (0;5] \)}{\(x\in (0;\infty )\)}{\(x\in (-\infty ;5] \)}{\(x\in [ 5;\infty )\)}{}{}}
\LANGpl{
\Question{
Rozwiąż następującą nierówność.
\[ 
 \log _{0.3}x\geq \log _{0.3}5
\]}
{\(x\in (0;5] \)}{\(x\in (0;\infty )\)}{\(x\in (-\infty ;5] \)}{\(x\in [ 5;\infty )\)}{}{}}
\LANGcs{
\Question{
Najděte všechna \(x\in \mathbb{R}\),
pro která platí \(\log _{0{,}3}x\geq \log _{0{,}3}5\).}
{\(x\in (0;5\rangle \)}{\(x\in (0;\infty )\)}{\(x\in (-\infty ;5\rangle \)}{\(x\in \langle 5;\infty )\)}{}{}}
\LANGsk{
\Question{
Množina všetkých riešení nerovnice 
\(\log _{0{,}3}x\geq \log _{0{,}3}5\) je:}
{\(x\in (0;5\rangle \)}{\(x\in (0;\infty )\)}{\(x\in (-\infty ;5\rangle \)}{\(x\in \langle 5;\infty )\)}{}{}}
\LANGes{
\Question{
Resuelve la siguiente inecuación.  
\[ 
 \log _{0.3}x\geq \log _{0.3}5
\]}
{\(x\in (0;5] \)}{\(x\in (0;\infty )\)}{\(x\in (-\infty ;5] \)}{\(x\in [ 5;\infty )\)}{}{}}
\End
